\section{AI 辅助开发}

本项目把 AI 工具作为工程协作环境的一部分，而不是把它作为算法结论的来源。量子线路
路由同时涉及量子信息、图搜索、学习模型、前后端接口和部署流程，单靠一次性编写代码
很难把这些部分稳定连接起来。AI 的作用主要体现在信息整理、方案比较、代码阅读、界面
修订和文档组织上；项目目标、算法取舍、实验口径和最终文字由作者根据课程要求和实际
运行结果确定。

\subsection{资料整理}

项目早期需要把课程知识、经典路由算法和工程实现放在同一问题框架中。AI 辅助完成的
第一类工作是资料整理：把量子态、酉门、双比特门、SWAP、硬件耦合图和线路深度这些
课程概念对应到路由编译任务；把 SABRE 的前沿门、候选交换和距离启发式整理成可复现
的基线描述；再把 NPQR 的初始映射、束搜索、动作评分和轨迹复放放到“学习评分搜索”
这一叙事下。

这种整理不是替代阅读论文，而是帮助形成报告结构。SABRE、图神经网络辅助路由和强化
学习路由等相关方法提供了参考坐标：传统启发式强调速度和稳定性，学习方法强调从局部
结构中学习动作偏好，搜索框架负责保证拓扑约束和状态推进。本项目最终选择的写法是先
复现 SABRE，再说明自研方法如何在同一实验口径下逐步改进。

\subsection{网页开发}

网页演示需要让读者直观看到量子路由的过程，而不是只看到一个数字表格。AI 在前端部分
主要用于交互方案和界面细节的讨论：输入区如何容纳 OpenQASM，拓扑图和指标面板如何
并列显示，轨迹回放如何同步显示 SWAP、物理边和逻辑位交换，错误信息如何写得清楚但
不干扰主流程。

前端技术栈保持简单：静态页面负责展示和交互，浏览器脚本负责请求服务接口、渲染拓扑、
回放轨迹和更新指标。这样的选择适合量子信息基础大作业，因为它不要求额外的前端构建流程，并且
可以直接通过 GitHub Pages 发布。AI 辅助的价值主要体现在快速比较布局方案、核对按钮状态、
梳理响应字段和改写界面文案；最终页面是否可用，仍由实际浏览器打开、接口调用和样例
回放决定。

\subsection{服务开发}

后端服务承担真实编译任务。AI 在这一部分更多地用于代码阅读和接口一致性梳理：
梳理 REST、MCP 和命令行三种调用方式需要共享哪些字段，整理状态查询、示例列表、线路校验、
同步编译、异步任务和拓扑查询之间的返回格式是否一致，帮助发现前端字段和后端响应之间
的不匹配。

服务端技术栈包括 Python、FastAPI、Uvicorn、Qiskit、PyTorch、NetworkX 和 Rustworkx。
其中 Qiskit 用于量子线路处理和 SABRE 基线，图工具用于硬件拓扑和距离计算，PyTorch
用于模型评分模块，FastAPI 和 Uvicorn 用于把编译能力暴露为服务。AI 辅助在这里更像
一名工程搭档：它可以根据接口约定整理验证项，但每一次接口修改都需要通过实际请求、
测试用例和网页调用来确认。

\subsection{部署测试}

部署阶段容易出现两类问题：一类是本地可运行但服务器环境缺少依赖，另一类是 README、
网页、接口和测试之间的说明不一致。AI 在这一阶段用于整理验证流程，例如服务健康状态、
样例编译、静态页面契约、公开测试和材料完整性。容器化服务、GitHub Pages 和公开仓库
共同构成了本项目的交付链路。

这一过程也帮助项目把“能运行”扩展为“能交付”。本地命令行适合快速复现实验，REST
服务适合网页和脚本调用，MCP 适合工具客户端集成，GitHub Pages 适合课程展示。AI 可以
辅助列出这些调用方式之间需要保持一致的结果字段，但真正的通过标准是服务能够启动、网页
能够调用、测试能够运行、报告中的实验数据能够对应到实际输出。

\subsection{报告整理}

报告写作阶段，AI 主要用于结构重组和文字修订。早期材料更像项目说明，包含较多接口、
路径和过程记录；正式报告需要改成量子信息基础大作业文风：先讲量子信息基础和硬件约束，再讲
SABRE 复现、自研方法、系统部署和实验对照。AI 在这里帮助发现章节之间的断裂，例如
摘要是否回应课程要求，图表标题是否过长，实验结论是否对应表格，系统与部署是否混在
同一节中。

最终写法遵循两个原则。第一，正文解释问题和方法，具体复现命令放到附录；第二，图表
标题保持简短，详细解释放在段落中。这样可以减少工程文件列表对正文节奏的干扰，让报告
更接近课程论文，而不是开发日志。

\subsection{过程控制}

AI 协作最关键的部分是控制权。作者负责确定量子信息基础大作业的目标、选定对照基线、决定哪些
实验进入正文、运行测试和编译 PDF，并对最终结论负责。AI 给出的代码方案、文字方案和
实验组织方案只有在符合实际代码、样例输出和课程要求时才会被采用。

因此，本项目中的 AI 使用方式可以概括为：用 AI 提高阅读、开发和整理效率，用代码运行
结果决定结论。对于量子线路路由这样涉及算法与工程闭环的问题，这种协作方式能够减少
重复劳动，也能帮助报告保持清晰结构；但项目可信度仍来自可运行系统、可复现实验和明确
的课程知识对应关系。
